\documentclass[10pt,a4paper]{article}

\usepackage[T1]{fontenc}
\usepackage{lmodern}
\usepackage[english]{babel}
\usepackage[a4paper,top=16mm,bottom=15mm,left=20mm,right=20mm]{geometry}
\usepackage{microtype}
\usepackage{amsmath}
\usepackage{booktabs}
\usepackage{tabularx}
\usepackage{array}
\usepackage{enumitem}
\usepackage{graphicx}
\usepackage{multicol}
\usepackage{caption}
\usepackage[hidelinks]{hyperref}
\usepackage{titlesec}

\setlength{\parindent}{0pt}
\setlength{\parskip}{3pt}
\setlist[itemize]{leftmargin=5mm,itemsep=1pt,topsep=1pt}
\renewcommand{\arraystretch}{1.06}
\captionsetup{font=small,labelfont=bf,hypcap=false}
\setlength{\intextsep}{5pt}
\setlength{\textfloatsep}{6pt}
\raggedbottom
\titleformat{\section}{\large\bfseries}{\thesection}{0.65em}{}
\titlespacing*{\section}{0pt}{7pt}{2pt}
\pagestyle{plain}
\newcolumntype{Y}{>{\raggedright\arraybackslash}X}

\begin{document}

\begin{center}
  {\LARGE\bfseries Q-learning in one maze\par}
  \vspace{5pt}
  {\normalsize Name: Lukasz Baldyga \rule{34mm}{0.4pt}\hfill
  Student ID: B24D-5026 \rule{24mm}{0.4pt}\par}
\end{center}

\section{Implementation:}

The Q-learning is implemented in \textbf{Rust} language with some Python libraries for visualization.

The code is in: https://github.com/dashingcontrol/q-learning-rust

The main file with the algorithm (written as a learning practice for myself in rust) is in main.rs file.

\section{Introduction}

The agent learns a route through one $6\times6$ maze. An open cell is a state. The available actions are up,
down, left, and right, but the program selects only moves that lead to another open cell. The start is
$S=(5,0)$ and the goal is $G=(0,5)$. Each move receives reward $-1$, so a shorter successful route has a
larger return.

Rust handles training and CSV export. Python only reads the saved data and draws the figures. Following slide
40, the program records Q-values before training, after episode 50, and after episode 100. It also saves the
training paths from episodes 1, 10, and 100.

\begin{center}
\captionof{table}{Parameters and the reason for using them in this example.}
\footnotesize
\begin{tabularx}{0.96\textwidth}{@{}p{0.23\textwidth}p{0.18\textwidth}Y@{}}
  \toprule
  \textbf{Setting} & \textbf{Value} & \textbf{Reason} \\
  \midrule
  Episodes & 100 & The assignment asks for the 100th training path. \\
  Learning rate & $\alpha=0.40$ & Changes the table within a short run while retaining part of the old estimate. \\
  Discount & $\gamma=0.90$ & Passes future step costs back through the maze. \\
  Exploration & $\max(0.40\cdot0.97^{e-1},0.02)$ & Explores more at the start and becomes more consistent later. \\
  Reward & $-1$ per legal move & Gives shorter successful routes a better return. \\
  Episode limit & 200 steps & Stops an unsuccessful episode from running forever. \\
  Random seed & 7 & Makes the example repeatable. \\
  \bottomrule
\end{tabularx}
\end{center}

\section{Q-learning method}

For a transition from $s$ to $s'$ after action $a$, the code applies the off-policy Q-learning update from
the lecture:
\begin{equation}
Q(s,a)\leftarrow Q(s,a)+\alpha\left[r+\gamma\max_{a'\in A(s')}Q(s',a')-Q(s,a)\right].
\end{equation}
Only permitted actions are included in $A(s')$. At the goal, the future term is zero. For the first update, all
stored values are zero, so
\[
Q_{\mathrm{new}}=0+0.40[-1+0.90(0)-0]=-0.40.
\]

Training uses epsilon-greedy action selection. Epsilon begins at $0.40$, is multiplied by $0.97$ after each
episode, and does not fall below $0.02$:
\[
\varepsilon_e=\max(0.40\cdot0.97^{e-1},0.02).
\]
Most moves therefore use a current best action, while some moves explore. Equal best actions are chosen
randomly during training. The final policy check sets $\varepsilon=0$.

\newpage
\section{Training results}

The first episode takes 36 steps. Episode 10 takes 22, and episode 100 takes 10. The first ten episodes
average 51.6 steps; the last ten average 10.2. All 100 training episodes reach the goal, and the first
10-step training route appears at episode 30.

\begin{center}
  \includegraphics[width=0.62\textwidth]{output/figures/01_learning_curve.png}
  \captionof{figure}{The light line shows every episode. The darker line averages the latest ten episodes.}
\end{center}

\begin{center}
  \includegraphics[width=0.94\textwidth]{output/figures/02_training_trajectories.png}
  \captionof{figure}{Training paths saved at episodes 1, 10, and 100. Repeated segments come from exploration.}
\end{center}

Figure 2 shows how the route changes during learning. In episode 1 the agent revisits several cells and needs
36 steps to reach the goal. By episode 10 it follows most of the useful route, but exploratory moves create
detours and the path still takes 22 steps. Episode 100 reaches the goal in 10 steps, which is the shortest
possible path length for this maze.

These paths include the exploration used during training. The separate evaluation on page 3 sets epsilon to
zero and follows only the best actions stored in the final Q-table.

\newpage
\section{Q-values and final policy}

\begin{center}
  \includegraphics[width=0.96\textwidth]{output/figures/03_q_values.png}
  \captionof{figure}{Best legal action and its Q-value before training, after episode 50, and after episode 100.}
\end{center}

Before training, every entry is zero. Values near the goal are learned first, then the estimates move back
through the maze. After episode 100, the best value at the start is $-5.6128$. The Q-value is not simply
$-10$ because future rewards are discounted by $\gamma=0.90$.

\begin{center}
\begin{minipage}[c]{0.46\textwidth}
  \centering
  \includegraphics[width=0.88\linewidth]{output/figures/04_greedy_evaluation.png}
  \captionof{figure}{Greedy evaluation after training.}
\end{minipage}\hfill
\begin{minipage}[c]{0.48\textwidth}
  \small
  \textbf{Final check}

  The greedy policy reaches the goal in 10 steps. It is therefore a shortest path for this layout. No
  exploration is used and the Q-table is left unchanged.
\end{minipage}
\end{center}

\section{Conclusion}

This report implements Q-learning loop in one small maze in Rust language. Epsilon-greedy exploration is
used during training. The separate check with epsilon set to zero then tests the learned policy without
random actions or further Q updates.

\setlength{\multicolsep}{3pt}
\begin{multicols}{2}[\section*{References and reproduction}]
\small
\textbf{Reference.} Y. Kobayashi, \textit{Reinforcement Learning}, lecture slides, 11 August 2026,
\columnbreak
\end{multicols}

\end{document}
